\begin{definition}[The inclusion $L^2(\gamma_n) \hookrightarrow L^1(\gamma_n)$]
Let $\gamma_n$ be as in \noderef{GaussMeasure}; it is a probability measure by
\noderef{GaussIsProb}, hence finite.  Consequently every $v \in L^2(\gamma_n)$ is
also $\gamma_n$-integrable, and we write $\iota v \in L^1(\gamma_n)$ for the class
of $v$ in $L^1(\gamma_n)$.  The map $\iota : L^2(\gamma_n) \to L^1(\gamma_n)$ is
the identity on representatives: $\iota v$ and $v$ are the same function, viewed
in the two different Lebesgue spaces.  It is linear and injective, and it is the
map used whenever an $L^2(\gamma_n)$ input is regarded as an $L^1(\gamma_n)$ input,
in particular in the statement that the Riesz transform extends from
$L^2(\gamma_n)$ to $L^1(\gamma_n)$.
\end{definition}
